\section*{Capítulo \ref{ch:g12:polymers} --- Os polímeros}

\begin{solution}{exo:g12:polymers:1}
Propeno, \ce{CH2=CH-CH3}.
\end{solution}

\begin{solution}{exo:g12:polymers:2}
\ce{-CF2-CF2-}; um \omterm{def:g12:polymers:addition-polymer}{polímero de adição}: o \omterm{def:g12:polymers:polymer}{monômero} tem uma ligação
\ce{C=C} e não se forma nenhum outro produto.
\end{solution}

\begin{solution}{exo:g12:polymers:3}
Adição: polietileno, PVC, poliestireno. Condensação: PET (um poliéster),
náilon (uma poliamida).
\end{solution}

\begin{solution}{exo:g12:polymers:4}
Celulose, amido, borracha de látex, proteínas. O náilon e o PVC são
sintéticos.
\end{solution}

\begin{solution}{exo:g12:polymers:5}
Cada ligação usa um grupo de cada \omterm{def:g12:polymers:polymer}{monômero}; com dois grupos, um \omterm{def:g12:polymers:polymer}{monômero}
pode se ligar dos dois lados, e a cadeia continua crescendo. Um \omterm{def:g12:polymers:polymer}{monômero}
com um só grupo interromperia a cadeia.
\end{solution}

\begin{solution}{exo:g12:polymers:6}
$M(\ce{C2H3Cl}) = 24.0 + 3.0 + 35.5 = \qty{62.5}{g/mol}$;
$n = 125000/62.5 = 2000$.
\end{solution}

\begin{solution}{exo:g12:polymers:7}
$M(\ce{C8H8}) = \qty{104.0}{g/mol}$; $M = 2500 \times 104.0 =
\qty{2.60e5}{g/mol}$.
\end{solution}

\begin{solution}{exo:g12:polymers:8}
O grupo ácido de uma \omterm{def:g7:atoms-and-molecules:molecule}{molécula} forma um \omterm{def:g11:functional-groups:carboxylic-acid}{éster} com o grupo \omterm{def:g11:functional-groups:alcohol}{álcool} da
seguinte, que ainda tem um grupo ácido livre: a cadeia cresce. \omterm{def:g12:polymers:repeat-unit}{Unidade
de repetição} \ce{-O-CH(CH3)-CO-}; libera-se água a cada ligação.
\end{solution}

\begin{solution}{exo:g12:polymers:9}
Os filmes de sacola são feitos de cadeias ramificadas, que não se
arrumam: o material é amorfo, macio e transparente. As garrafas de leite
são feitas de cadeias quase lineares, que se arrumam em regiões
cristalinas: mais densas, mais rígidas, e as regiões cristalinas
espalham a luz.
\end{solution}

\begin{solution}{exo:g12:polymers:10}
As suas cadeias estão unidas por muitas ligações cruzadas covalentes. O
aquecimento não consegue fazer as cadeias deslizarem; ele rompe ligações
e destrói o material em vez de fundi-lo.
\end{solution}

\begin{solution}{exo:g12:polymers:11}
$M(\ce{C12H22N2O2}) = 144.0 + 22.0 + 28.0 + 32.0 = \qty{226.0}{g/mol}$;
$n = 22600/226.0 = 100$.
\end{solution}

\begin{solution}{exo:g12:polymers:12}
\[
  \ce{H2N(CH2)6NH2 + HOOC(CH2)4COOH -> H2N(CH2)6NHCO(CH2)4COOH + H2O} .
\]
O produto ainda carrega um grupo \omterm{def:g11:functional-groups:amine}{amina} livre numa ponta e um grupo
ácido livre na outra: cada um pode reagir com outro \omterm{def:g12:polymers:polymer}{monômero}.
\end{solution}

\begin{solution}{exo:g12:polymers:13}
Duas \omterm{def:g7:atoms-and-molecules:molecule}{moléculas} de água por \omterm{def:g12:polymers:repeat-unit}{unidade de repetição}. $n = 1000/226.0 =
\qty{4.42}{mol}$ de \omterm{def:g12:polymers:repeat-unit}{unidades de repetição}, logo \qty{8.85}{mol} de água:
$8.85 \times 18.0 = \qty{159}{g}$.
\end{solution}

\begin{solution}{exo:g12:polymers:14}
A borracha natural quente é feita de cadeias que deslizam umas sobre as
outras. As pontes de enxofre ligam as cadeias: elas se esticam mas
voltam, e não escorrem mais. Como tem ligações cruzadas, um pneu não
pode ser fundido e remodelado.
\end{solution}

\begin{solution}{exo:g12:polymers:15}
$M(\ce{C6H10O5}) = \qty{162.0}{g/mol}$; $n = \num{1.6e6}/162.0 = 9900$
unidades; comprimento $9900 \times 0.5 = \qty{4900}{nm}$, cerca de
\qty{5}{\micro m}.
\end{solution}

\begin{solution}{pb:g12:polymers:1}
\textbf{1.} Etano-1,2-diol (dois grupos \omterm{def:g11:functional-groups:alcohol}{álcool}) e ácido
benzeno-1,4-dicarboxílico (dois grupos \omterm{def:g11:functional-groups:carboxylic-acid}{ácido carboxílico}).

\textbf{2.} Uma ligação \omterm{def:g11:functional-groups:carboxylic-acid}{éster}; libera-se água.

\textbf{3.} Um \omterm{def:g12:polymers:addition-polymer}{polímero de condensação}.

\textbf{4.} $M(\ce{C2H6O2}) = \qty{62.0}{g/mol}$; $M(\ce{C8H6O4}) =
\qty{166.0}{g/mol}$.

\textbf{5.} $62.0 + 166.0 - 2 \times 18.0 = \qty{192.0}{g/mol}$; e
$10 \times 12.0 + 8 \times 1.0 + 4 \times 16.0 = 192.0$.

\textbf{6.} $n = 38400/192.0 = 200$.

\textbf{7.} Duas ligações \omterm{def:g11:functional-groups:carboxylic-acid}{éster} por \omterm{def:g12:polymers:repeat-unit}{unidade de repetição}: cerca de 400.

\textbf{8.} As suas cadeias são em grande parte lineares, e partes delas
se alinham em regiões ordenadas.

\textbf{9.} Ele é um \omterm{def:g8:plastics:thermoplastic}{termoplástico}: as suas cadeias não têm ligações
cruzadas e deslizam umas sobre as outras quando quentes.

\textbf{10.} $300/0.80 = \qty{375}{g}$.

\textbf{11.} $375/25 = 15$ garrafas.

\textbf{12.} Tampas, rótulos, sujeira e aparas são retirados ou perdidos
durante a triagem, a lavagem e a fiação.

\textbf{13.} Prevenir os resíduos (princípio 1): um objeto usado vira
\omterm{def:g5:raw-materials-and-recycling:raw-material}{matéria-prima}.

\textbf{14.} \ce{C10H8O4 + 2H2O -> C2H6O2 + C8H6O4}.

\textbf{15.} $1000/192.0 = \qty{5.21}{mol}$.

\textbf{16.} Diol: $5.21 \times 62.0 = \qty{323}{g}$; ácido:
$5.21 \times 166.0 = \qty{865}{g}$.

\textbf{17.} $2 \times 5.21 \times 18.0 = \qty{188}{g}$ de água;
$1000 + 188 = 323 + 865 = \qty{1188}{g}$: a massa se conserva.

\textbf{18.} Os \omterm{def:g12:polymers:polymer}{monômeros} podem ser purificados e dão um PET novo tão
bom quanto o original, enquanto a fusão encurta um pouco as cadeias a
cada vez.

\textbf{19.} Uma jaqueta de fleece leva 15 garrafas de \qty{25}{g}.
\end{solution}
